\subsubsection{QUBO Problems for QAOA}\label{sec:qubo-problems-for-qaoa}

This section is where QAOA reconnects with everything we saw in QUBO Formulation (\autopageref{sec:qubo-formulation}). In fact, there is almost no new mathematics here, but rather we answer an important practical question: \emph{\textbf{which optimization problems can we naturally solve with QAOA?}} The answer is: \textbf{QUBO problems}.

\highspace
\textcolor{Green3}{\faIcon{question-circle} \textbf{Why QUBO?}} At the beginning of the QAOA section, we saw that given an optimization problem, we can construct a cost Hamiltonian $H_C$ and a mixer Hamiltonian $H_M$, run the QAOA circuit, and measure the output state to get a solution. But there is one missing step. \emph{\textbf{How do we build the Cost Hamiltonian?}} The answer is simple: given an optimization problem, we can always write it in QUBO form, and then we can construct the Cost Hamiltonian from the QUBO matrix $Q$. Therefore, \textbf{QUBO} is \textbf{not the algorithm}, but rather the \textbf{standard mathematical representation} used before constructing the Cost Hamiltonian.

\highspace
\textcolor{Green3}{\faIcon{question-circle} \textbf{But why QUBO and not some other representation?}} In this course, we will focus on QUBO because it is well suited for QAOA. The reason is that \hl{QUBO already has exactly the mathematical structure that we need to construct a Hamiltonian.} Let's see why.

\highspace
\textcolor{Green3}{\faIcon{book} \textbf{The general QUBO formulation.}} Recall that a QUBO problem (\autopageref{sec:introduction-to-qubo}) is defined as follows:
\begin{equation*}
    \min_{x\in\{0,1\}^n} x^T Q x
\end{equation*}
Where $x$ is a binary vector $x \in \left\{0,1\right\}^n$ and $Q$ is an $n \times n$ matrix containing the linear and quadratic coefficients of the objective function.\footnote{%
    The matrix $Q$ contains all the coefficients describing the optimization problem.
} Also, we know that the main diagonal entries $Q_{ii}$ represent the \textbf{linear terms} $x_i$ (i.e., contribution of individual variables), while the off-diagonal entries $Q_{ij}$ represent the \textbf{quadratic terms} $x_i x_j$ (i.e., interaction between pairs of variables).

\begin{examplebox}[: QUBO representation of a polynomial]
    Suppose our object function is:
    \begin{equation*}
        E(x) = x_1 + x_2 - 2x_1 x_2
    \end{equation*}
    This polynomial is the optimization problem. Now, let's encode it into a matrix. We want to find a way to represent:
    \begin{equation*}
        E(x) = x^T Q x = x_1 + x_2 - 2x_1 x_2 \quad \text{with } x = \begin{bmatrix} x_1 \\ x_2 \end{bmatrix}
    \end{equation*}
    There are \textbf{two} ways to write this polynomial as a quadratic form.
    \begin{enumerate}
        \item \textbf{Symmetric matrix}. Substitute the values of $x_1$ and $x_2$ into the polynomial and rewrite it as a quadratic form:
        \begin{equation*}
            E(x) = x^T \begin{bmatrix} 1 & -1 \\ -1 & 1 \end{bmatrix} x
        \end{equation*}
        Where the matrix $Q$ is composed as follows:
        \begin{itemize}
            \item The diagonal entries $Q_{11} = 1$ and $Q_{22} = 1$ correspond to the linear terms $x_1$ and $x_2$.
            \item The off-diagonal entries $Q_{12} = -1$ and $Q_{21} = -1$ correspond to the quadratic term $-2x_1 x_2$.
        \end{itemize}
        Now compote $Qx$:
        \begin{equation*}
            Qx = \begin{bmatrix} 1 & -1 \\ -1 & 1 \end{bmatrix} \begin{bmatrix} x_1 \\ x_2 \end{bmatrix} = \begin{bmatrix} x_1 - x_2 \\ -x_1 + x_2 \end{bmatrix}
        \end{equation*}
        Then compute $x^T Q x$:
        \begin{equation*}
            \begin{aligned}
                x^T Q x &= \begin{bmatrix} x_1 & x_2 \end{bmatrix} \begin{bmatrix} x_1 - x_2 \\ -x_1 + x_2 \end{bmatrix} \\
                &= x_1\left(x_1 - x_2\right) + x_2\left(-x_1 + x_2\right) \\
                &= x_1^2 - x_1 x_2 - x_2 x_1 + x_2^2 \\
                &= x_1 + x_2 - 2x_1 x_2
            \end{aligned}
        \end{equation*}
        So we have successfully represented the polynomial as a quadratic form using a symmetric matrix. However, notice something: the interaction appeared \textbf{twice} $-x_1 x_2$ and $-x_2 x_1$, is stored in two different entries of the matrix.


        \item \textbf{Upper triangular}. Instead, we can write the polynomial as a quadratic form using an upper triangular matrix:
        \begin{equation*}
            E(x) = x^T \begin{bmatrix} 1 & -2 \\ 0 & 1 \end{bmatrix} x
        \end{equation*}
        Where the matrix $Q$ is composed as follows:
        \begin{itemize}
            \item The diagonal entries $Q_{11} = 1$ and $Q_{22} = 1$ correspond to the linear terms $x_1$ and $x_2$.
            \item The off-diagonal entry $Q_{12} = -2$ corresponds to the quadratic term $-2x_1 x_2$.
        \end{itemize}
        Now compote $Qx$:
        \begin{equation*}
            Qx = \begin{bmatrix} 1 & -2 \\ 0 & 1 \end{bmatrix} \begin{bmatrix} x_1 \\ x_2 \end{bmatrix} = \begin{bmatrix} x_1 - 2x_2 \\ x_2 \end{bmatrix}
        \end{equation*}
        Then compute $x^T Q x$:
        \begin{equation*}
            \begin{aligned}
                x^T Q x &= \begin{bmatrix} x_1 & x_2 \end{bmatrix} \begin{bmatrix} x_1 - 2x_2 \\ x_2 \end{bmatrix} \\
                &= x_1\left(x_1 - 2x_2\right) + x_2\left(x_2\right) \\
                &= x_1^2 - 2x_1 x_2 + x_2^2 \\
                &= x_1 + x_2 - 2x_1 x_2
            \end{aligned}
        \end{equation*}
        So we have successfully represented the polynomial as a quadratic form using an upper triangular matrix. Notice that the interaction appeared \textbf{once} $-2x_1 x_2$, and is stored in a single entry of the matrix.
    \end{enumerate}
    So when we construct $Q$, we have to decide \textbf{how to store the interaction coefficients}: symmetric (split the coefficient equally between $Q_{ij}$ and $Q_{ji}$) or upper triangular (store the coefficient in a single entry $Q_{ij}$). The choice is arbitrary.
\end{examplebox}

\highspace
In summary, QUBO provides a standard mathematical formulation for many combinatorial optimization problems. A QUBO problem is written as:
\begin{equation*}
    \min_{x \in \{0,1\}^n} x^T Q x
\end{equation*}
Where $x$ is a binary vector and $Q$ contains the linear and quadratic coefficients of the objective function. Since this formulation can be systematically converted into spin variables and then into a Cost Hamiltonian, QUBO is particularly well suited for QAOA. The matrix $Q$ may be symmetric or upper triangular because the interaction terms $x_i x_j$ are symmetric.